%------------------------------------------------------------------------------
% 4.8 مصفوفة تتبع التصميم إلى المتطلبات
% المالك: TM (مهمة جانت ٢.٢.٧ — المراجعة والاعتماد 12/11 إلى 14/11/2026)
% الغرض: القفل النهائي للفصل — إثبات قابل للفحص أن كل متطلب وظيفي
%   (FR-01..FR-29) وكل حالة استخدام (UC-01..UC-07) لهما موضع تصميم صريح
%   في هذا الفصل، فلا ينتقل شيء إلى التنفيذ بلا أساس تصميمي.
%
% عناصر مطلوبة:
%   [ ] الجدول المرجعي (tab:design-traceability): المتطلب / حالة الاستخدام ←
%       عنصر التصميم المخدوم (رقم القسم أو الجدول أو الشكل) ← ملاحظة —
%       29 صفاً للمتطلبات و7 صفوف لحالات الاستخدام (يجوز تقسيمه جدولين
%       للمتطلبات الوظيفية وحالات الاستخدام)
%   [ ] فقرة تحليلية قصيرة: قراءة المصفوفة — أي متطلبات تخدمها عناصر
%       تصميم متعددة (قوة)، وأيها يخدمها عنصر واحد (نقطة انتباه في
%       الاختبار)، وأي متطلب بلا موضع تصميم (فجوة يجب سدها قبل الفصل
%       الخامس — إن وُجد)
%
% أسلوب الجدول: نفس أسلوب «جدول الملخص المرجعي السريع» الذي نجح في
%   حالات الاستخدام (3.4.4) — أعمدة قصيرة وصف لكل عنصر
%------------------------------------------------------------------------------
\section{مصفوفة تتبع التصميم إلى المتطلبات}
\label{sec:design-traceability}
~
